\section*{Hoofdstuk \ref{ch:g10:why-clouds-form} --- Waarom wolken ontstaan}

\begin{solution}{exo:g10:why-clouds-form:1}
De \omterm{def:g10:why-clouds-form:humidity}{relatieve luchtvochtigheid} is het gehalte aan waterdamp van de lucht gedeeld
door haar verzadigingsgehalte bij dezelfde temperatuur en druk, als percentage.
Het \omterm{def:g10:why-clouds-form:dew-point}{dauwpunt} is de temperatuur waartoe de lucht moet worden afgekoeld, met behoud
van haar watergehalte en druk, om \omterm{def:g10:why-clouds-form:humidity}{verzadigd} te raken.
\end{solution}

\begin{solution}{exo:g10:why-clouds-form:2}
\num{3.8}, \num{7.6}, \num{14.7} en \qty{27.2}{g/kg}. Elke
\qty{10}{\degreeCelsius} vermenigvuldigt het met ongeveer 2 ($7.6/3.8 = 2.0$,
$14.7/7.6 \approx 1.9$, $27.2/14.7 \approx 1.9$).
\end{solution}

\begin{solution}{exo:g10:why-clouds-form:3}
Verzadigingsgehalte bij \qty{24}{\degreeCelsius}: \qty{18.9}{g/kg}; \omterm{def:g10:why-clouds-form:humidity}{relatieve
luchtvochtigheid} $10.0/18.9 \approx \qty{53}{\percent}$. Het verzadigingsgehalte
is \qty{10.0}{g/kg} bij \qty{14}{\degreeCelsius}: dat is het \omterm{def:g10:why-clouds-form:dew-point}{dauwpunt}.
\end{solution}

\begin{solution}{exo:g10:why-clouds-form:4}
In de zomer is de lucht vochtig en ligt haar \omterm{def:g10:why-clouds-form:dew-point}{dauwpunt} ruim boven de temperatuur
van het glas: de lucht die het glas raakt, wordt onder haar \omterm{def:g10:why-clouds-form:dew-point}{dauwpunt} afgekoeld en
er condenseert water op. Heel droge winterlucht binnenshuis heeft een \omterm{def:g10:why-clouds-form:dew-point}{dauwpunt}
dat lager is dan de temperatuur van het glas: ook als ze ertegen afkoelt, is de
lucht nog niet \omterm{def:g10:why-clouds-form:humidity}{verzadigd}.
\end{solution}

\begin{solution}{exo:g10:why-clouds-form:5}
Verwarming van onderen (middagcumulus boven zonnige grond); bergen (wind die
langs een bergketen omhoog wordt gedwongen); fronten (warme lucht die boven koude
lucht uitstijgt); convergentie (winden die bij de evenaar of in een depressie
samenstromen).
\end{solution}

\begin{solution}{exo:g10:why-clouds-form:6}
$125 \times (30 - 18) = \qty{1500}{m}$; $125 \times (22 - 20) =
\qty{250}{m}$. De tweede dag: de wolken ontstaan maar \qty{250}{m} boven de
grond.
\end{solution}

\begin{solution}{exo:g10:why-clouds-form:7}
Op \qty{1}{km}: $18 - 9.8 = \qty{8.2}{\degreeCelsius}$; op \qty{2}{km}:
$18 - 19.6 = \qty{-1.6}{\degreeCelsius}$. Bij het dalen warmt het in hetzelfde
tempo op en bereikt het de grond weer bij \qty{18}{\degreeCelsius}.
\end{solution}

\begin{solution}{exo:g10:why-clouds-form:8}
Bij \qty{4}{\degreeCelsius/km} koelen zowel een onverzadigd pakketje
(\qty{9.8}{\degreeCelsius/km}) als een \omterm{def:g10:why-clouds-form:humidity}{verzadigd} pakketje
(\qty{5}{\degreeCelsius/km}) sneller af dan hun omgeving: in beide gevallen
\omterm{def:g10:why-clouds-form:stability}{stabiel}. Bij \qty{12}{\degreeCelsius/km} koelt zelfs een onverzadigde bel
langzamer af dan haar omgeving: ze blijft warmer en blijft stijgen, en de bellen
warme lucht stijgen vrij op.
\end{solution}

\begin{solution}{exo:g10:why-clouds-form:9}
$14.7 - 7.6 = \qty{7.1}{g}$ per kilogram lucht. De \omterm{def:g10:earth-energy-budget:latent-heat}{latente warmte} verwarmt de
lucht waarin het water condenseert, zodat die trager afkoelt.
\end{solution}

\begin{solution}{exo:g10:why-clouds-form:10}
Volumes groeien met de derde macht van de diameter: $(0.5/0.02)^3 = 25^3 = 15\,625$
druppeltjes voor een motregendruppel; $(2/0.02)^3 = 100^3 = 1\,000\,000$ voor een
regendruppel.
\end{solution}

\begin{solution}{exo:g10:why-clouds-form:11}
Onder een heldere hemel straalt de grond infrarood uit en koelt af, en de lucht
die haar raakt ook. Die koude lucht is dichter dan de lucht erboven; op hellingen
stroomt ze naar beneden en verzamelt ze zich op de bodem van het dal. Daar koelt
ze de hele nacht verder af tot ze haar \omterm{def:g10:why-clouds-form:dew-point}{dauwpunt} bereikt, en ontstaan er
druppeltjes in de lucht zelf: mist. Op de heuveltoppen is de lucht niet zo koud
en blijft ze boven haar \omterm{def:g10:why-clouds-form:dew-point}{dauwpunt}.
\end{solution}

\begin{solution}{exo:g10:why-clouds-form:12}
Op hoogte $z$ is de temperatuur $T - a z$ en het \omterm{def:g10:why-clouds-form:dew-point}{dauwpunt} $T_d - b z$; die zijn
gelijk wanneer $(a - b) z = T - T_d$, dus $z = (T - T_d)/(a - b)$. Met
$b = 1.7$: $a - b = 8.1$, en $1000/8.1 \approx \qty{123}{m}$ per graad.
\end{solution}

\begin{solution}{exo:g10:why-clouds-form:13}
\omterm{def:g10:why-clouds-form:condensation-level}{Condensatieniveau}: $125 \times (15 - 11) = \qty{500}{m}$; temperatuur daar:
$15 - 9.8 \times 0.5 \approx \qty{10.1}{\degreeCelsius}$. Top: $10.1 - 5
\times 1.5 = \qty{2.6}{\degreeCelsius}$. \qty{2}{km} naar beneden met de droge
gradiënt: $2.6 + 9.8 \times 2 \approx \qty{22.2}{\degreeCelsius}$, ongeveer
zeven graden warmer dan aan de loefzijde.
\end{solution}

\begin{solution}{exo:g10:why-clouds-form:14}
Bij \qty{-20}{\degreeCelsius}: $126/103 \approx 1.22$, \qty{22}{\percent}
oververzadigd voor het ijs; bij \qty{-10}{\degreeCelsius}: $287/260 \approx
1.10$, \qty{10}{\percent}. Binnen dit temperatuurbereik geldt: hoe kouder de wolk,
hoe groter het voordeel van het ijs. Het proces werkt het best in het koude
bovenste deel van een wolk, waar druppeltjes en kristallen door elkaar zitten.
\end{solution}

\begin{solution}{exo:g10:why-clouds-form:15}
Onverzadigd: het pakketje koelt \qty{9.8}{\degreeCelsius} af, zijn omgeving
\qty{6.5}{\degreeCelsius}: het is \qty{3.3}{\degreeCelsius} kouder en zakt terug.
\omterm{def:g10:why-clouds-form:humidity}{Verzadigd}: het koelt \qty{5}{\degreeCelsius} af, dus het is
\qty{1.5}{\degreeCelsius} warmer en blijft stijgen. De lucht is alleen \omterm{def:g10:why-clouds-form:stability}{onstabiel}
op één voorwaarde: dat de stijgende lucht \omterm{def:g10:why-clouds-form:humidity}{verzadigd} is.
\end{solution}

\begin{solution}{pb:g10:why-clouds-form:1}
\textbf{1.} \qty{14.7}{g/kg}.

\textbf{2.} $11.4/14.7 \approx \qty{78}{\percent}$.

\textbf{3.} \qty{16}{\degreeCelsius}, waar het verzadigingsgehalte
\qty{11.4}{g/kg} is.

\textbf{4.} $11.4/24.1 \approx \qty{47}{\percent}$.

\textbf{5.} Nee: ze bevat dezelfde \qty{11.4}{g/kg}. Haar \omterm{def:g10:why-clouds-form:humidity}{relatieve
luchtvochtigheid} daalde omdat ze, warmer, meer kon bevatten.

\textbf{6.} Bij haar \omterm{def:g10:why-clouds-form:dew-point}{dauwpunt}, \qty{16}{\degreeCelsius}.

\textbf{7.} $11.4 - 8.7 = \qty{2.7}{g}$ per kilogram lucht.

\textbf{8.} Onder een heldere nachthemel stralen gras en daken infrarood uit en
worden ze kouder dan de lucht; de lucht die ze raakt, wordt onder haar \omterm{def:g10:why-clouds-form:dew-point}{dauwpunt}
afgekoeld voordat de rest van de lucht zover is.

\textbf{9.} $(28.0 - 13.4)/1.5 \approx \qty{9.7}{\degreeCelsius/km}$,
vrijwel de droge gradiënt: de middagbellen hebben deze laag gemengd.

\textbf{10.} $(13.4 - 0.4)/2.0 = \qty{6.5}{\degreeCelsius/km}$.

\textbf{11.} \omterm{def:g10:why-clouds-form:stability}{Stabiel} voor een onverzadigd pakketje, dat sneller afkoelt
(\qty{9.8}{\degreeCelsius/km}); \omterm{def:g10:why-clouds-form:stability}{onstabiel} voor een \omterm{def:g10:why-clouds-form:humidity}{verzadigd} pakketje, dat
langzamer afkoelt (\qty{5}{\degreeCelsius/km}).

\textbf{12.} De temperatuur stijgt met de hoogte, van \num{-0.3} tot
\qty{1.2}{\degreeCelsius}: een zeer stabiele laag, die als een deksel werkt. De
stijgende wolk stopt daar: haar top ligt bij ongeveer \qty{4}{km}.

\textbf{13.} $28 - 16 = \qty{12}{\degreeCelsius}$.

\textbf{14.} Temperatuur $28 - 9.8 = \qty{18.2}{\degreeCelsius}$; \omterm{def:g10:why-clouds-form:dew-point}{dauwpunt}
$16 - 1.8 = \qty{14.2}{\degreeCelsius}$.

\textbf{15.} Van 12 naar $18.2 - 14.2 = 4.0$: het verschil kromp met
\qty{8}{\degreeCelsius}, een snelheid van \qty{8}{\degreeCelsius} per kilometer.

\textbf{16.} Eén graad verschil wordt in $1/8$ km overbrugd, dat is
\qty{125}{m}.

\textbf{17.} $125 \times (20 - 16) = \qty{500}{m}$.

\textbf{18.} Het \omterm{def:g10:why-clouds-form:dew-point}{dauwpunt} blijft \qty{16}{\degreeCelsius} terwijl de lucht bij
de grond opwarmt: het verschil wordt groter, en daarmee ook de hoogte waarop het
wordt overbrugd.

\textbf{19.} $28 - 9.8\,z = 16 - 1.8\,z$.

\textbf{20.} $8\,z = 12$, dus $z = 1.5$: de basis van de middagcumulus ligt
\textbf{ongeveer \qty{1.5}{km}} boven de vlakte (de exacte berekening met de
eigenschappen van water geeft \qty{1.51}{km}).
\end{solution}
